
\documentclass{article}
\usepackage{colortbl}
\usepackage{makecell}
\usepackage{multirow}
\usepackage{supertabular}

\begin{document}

\newcounter{utterance}

\centering \large Interaction Transcript for game `function\_detective', experiment `string\_passive\_labeled\_pairs\_oneshot', episode 179 with mistral{-}large{-}3{-}azure.
\vspace{24pt}

{ \footnotesize  \setcounter{utterance}{1}
\setlength{\tabcolsep}{0pt}
\begin{supertabular}{c@{$\;$}|p{.15\linewidth}@{}p{.15\linewidth}p{.15\linewidth}p{.15\linewidth}p{.15\linewidth}p{.15\linewidth}}
    \# & \multicolumn{2}{c}{Player} && \multicolumn{2}{c}{Game Master} \\
    \hline

    \theutterance \stepcounter{utterance}  
    & & & \multicolumn{4}{p{0.6\linewidth}}{
        \cellcolor[rgb]{0.9,0.9,0.9}{
            \makecell[{{p{\linewidth}}}]{
                \texttt{\tiny{[P1$\langle$GM]}}
                \texttt{I have a hidden Python function for you to figure out.} \\
\\ 
\texttt{Your job is to infer what it does and then submit a correct Python implementation.} \\
\\ 
\texttt{Function signature:} \\
\texttt{(text: str) {-}> str} \\
\\ 
\texttt{You have 1 turns total.} \\
\\ 
\texttt{I have already given you a fixed list of candidate input{-}output pairs with true or false labels. Solve using those statements only.} \\
\\ 
\texttt{You must make your best SOLVE attempt no later than the second{-}to{-}last turn.} \\
\texttt{Do not wait until the very end.} \\
\\ 
\texttt{You can use these actions:} \\
\texttt{{-} JSON format described at the end of the prompt} \\
\\ 
\\ 
\texttt{Do not add import statements in your guess.} \\
\texttt{Anything you might need is already available, so just write the function body in hypothesis and return the answer.} \\
\\ 
\texttt{Preloaded examples:} \\
\texttt{1. (('hello world', '\#hello world'), True)} \\
\texttt{2. ((' ', '\# \_alt'), False)} \\
\texttt{3. (('A', '\#A'), True)} \\
\texttt{4. (('123', '\#123\_alt'), False)} \\
\texttt{5. (('aba', '\#aba'), True)} \\
\texttt{6. (('aA', '\#aA\_alt'), False)} \\
\texttt{7. (('', '\#'), True)} \\
\texttt{8. (('hello', '\#hello\_alt'), False)} \\
\texttt{9. (('a{-}b', '\#a{-}b'), True)} \\
\texttt{10. (('ab', '\#ab\_alt'), False)} \\
\\ 
\texttt{Scoring:} \\
\texttt{{-} correctness is required for any non{-}zero score} \\
\texttt{{-} in interactive modes, solving with fewer information requests is better} \\
\\ 
\texttt{Input formatting reminders:} \\
\texttt{{-} strings must be quoted} \\
\texttt{{-} booleans must be True/False} \\
\texttt{{-} use commas between multiple values} \\
\texttt{{-} do not write any extra action tags in the same message} \\
\\ 
\texttt{Format example:} \\
\texttt{You write:} \\
\texttt{\{'current\_mode': <probing/confirming/don't know; choose whichever best describes your current state in the task>, 'input\_rationale': <brief description of why you chose this hypothesis>, 'hypothesis': <either the exact string "unknown" if you have no guess yet OR your current best guess as runnable Python matching the given signature, with a one{-}line docstring stating what it does. Never write placeholder code, if you cannot commit to a guess, use "unknown". Match the given signature.} \\
\texttt{e.g.} \\
\texttt{def f(x):} \\
\texttt{	"""Returns the absolute value of x."""} \\
\texttt{	return x if x >= 0 else {-}x} \\
\\ 
\texttt{Use "unknown" if you have no guess.>, 'confidence\_score': <score between 1{-}100. 100 means highest confidence in your hypothesis>\}} \\
\\ 
\texttt{Output only one valid json per turn.} \\
            }
        }
    }
     \\ \\

    \theutterance \stepcounter{utterance}  
    & \multicolumn{4}{p{0.6\linewidth}}{
        \cellcolor[rgb]{0.561,0.737,0.561}{
            \makecell[{{p{\linewidth}}}]{
                \texttt{\tiny{[P1$\rangle$GM]}}
                \texttt{```json} \\
\texttt{\{} \\
\texttt{  "current\_mode": "confirming",} \\
\texttt{  "input\_rationale": "Based on the provided input{-}output pairs, the function appears to prepend a '\#' to the input string if the input string meets certain conditions. The conditions seem to be: (1) The input string is non{-}empty and all characters are alphanumeric or hyphens (no underscores or other special characters), or (2) the input string is empty. If the input string contains any characters that are not alphanumeric or hyphens, or if the output has an '\_alt' suffix, the function returns False. The function seems to validate that the output is exactly '\#' followed by the input string, with no additional suffixes or modifications.",} \\
\texttt{  "hypothesis": "def f(text: str) {-}> str:\textbackslash{}n    \textbackslash{}"\textbackslash{}"\textbackslash{}"Prepends a '\#' to the input string if it consists only of alphanumeric characters or hyphens, or if it is empty. Otherwise, the behavior is undefined based on the given examples (but the examples suggest it should not add '\_alt').\textbackslash{}"\textbackslash{}"\textbackslash{}"\textbackslash{}n    return '\#' + text",} \\
\texttt{  "confidence\_score": 90} \\
\texttt{\}} \\
\texttt{```} \\
            }
        }
    }
    & & \\ \\

    \theutterance \stepcounter{utterance}  
    & & & \multicolumn{2}{p{0.3\linewidth}}{
        \cellcolor[rgb]{0.9,0.9,0.9}{
            \makecell[{{p{\linewidth}}}]{
                \texttt{\tiny{[GM$|$GM]}}
                \texttt{Game Verdict: WIN} \\
            }
        }
    }
    & & \\ \\

    \theutterance \stepcounter{utterance}  
    & & & \multicolumn{2}{p{0.3\linewidth}}{
        \cellcolor[rgb]{0.9,0.9,0.9}{
            \makecell[{{p{\linewidth}}}]{
                \texttt{\tiny{[GM$|$GM]}}
                \texttt{That is correct!} \\
            }
        }
    }
    & & \\ \\

    \theutterance \stepcounter{utterance}  
    & & & \multicolumn{2}{p{0.3\linewidth}}{
        \cellcolor[rgb]{0.9,0.9,0.9}{
            \makecell[{{p{\linewidth}}}]{
                \texttt{\tiny{[GM$|$GM]}}
                \texttt{def prepend\_hash(text: str) {-}> str:} \\
\texttt{    return '\#' + text} \\
            }
        }
    }
    & & \\ \\

    \theutterance \stepcounter{utterance}  
    & & & \multicolumn{2}{p{0.3\linewidth}}{
        \cellcolor[rgb]{0.9,0.9,0.9}{
            \makecell[{{p{\linewidth}}}]{
                \texttt{\tiny{[GM$|$GM]}}
                \texttt{def f(text: str) {-}> str:} \\
\texttt{    """Prepends a '\#' to the input string if it consists only of alphanumeric characters or hyphens, or if it is empty. Otherwise, the behavior is undefined based on the given examples (but the examples suggest it should not add '\_alt')."""} \\
\texttt{    return '\#' + text} \\
            }
        }
    }
    & & \\ \\

    \theutterance \stepcounter{utterance}  
    & & & \multicolumn{2}{p{0.3\linewidth}}{
        \cellcolor[rgb]{0.9,0.9,0.9}{
            \makecell[{{p{\linewidth}}}]{
                \texttt{\tiny{[GM$|$GM]}}
                \texttt{Mode: passive\_labeled\_pairs\_oneshot} \\
\texttt{Accuracy: 1.000} \\
\texttt{Efficiency: 1.000} \\
\texttt{Parse errors: 0} \\
\texttt{Runtime errors: 0} \\
            }
        }
    }
    & & \\ \\

\end{supertabular}
}

\end{document}
